\documentclass[11pt]{article}
\usepackage[margin=1in]{geometry}
\usepackage{booktabs}
\usepackage{graphicx}
\usepackage{amsmath}
\usepackage{siunitx}
\usepackage{enumitem}

\title{Real-World UR10e Experimental Plan\\[0.3em]
\large RIT*: Riemannian Informed Trees}
\author{}
\date{}

\begin{document}
\maketitle

\section*{General Setup}

\begin{itemize}[itemsep=2pt]
  \item \textbf{Robot:} Universal Robots UR10e + Robotiq 85 gripper
  \item \textbf{Control:} \texttt{ur\_rtde} Python driver at \SI{500}{\hertz}
  \item \textbf{Coordinate frame:} Robot base at origin $[0,0,0]$. $X$ = forward, $Y$ = left, $Z$ = up
  \item \textbf{All positions in meters}, measured from the base center
  \item \textbf{Speed limit:} Max \SI{0.5}{\radian\per\second} per joint during initial tests; \SI{1.0}{\radian\per\second} for data collection
  \item \textbf{Planners compared:} RIT* (ours), BIT*, RRT*-Connect
  \item \textbf{Trials:} 20 random seeds per planner per experiment
  \item \textbf{Obstacle geometry:} Registered to the planner manually via tape-measure from the base center ($\pm$\SI{2}{\centi\meter} accuracy)
\end{itemize}

\subsection*{Common Metrics (all experiments)}

\begin{enumerate}[itemsep=2pt]
  \item \textbf{Planning time} (seconds) --- wall-clock from planner start to solution
  \item \textbf{Success rate} --- fraction of seeds that find a valid path within the iteration budget
  \item \textbf{Total joint torque} --- $\sum_{t} |\boldsymbol{\tau}(t) \cdot \Delta\mathbf{q}(t)|$ read from \texttt{ur\_rtde} at \SI{500}{\hertz}
  \item \textbf{Peak joint velocity} --- $\max_t \|\dot{\mathbf{q}}(t)\|_\infty$
  \item \textbf{Path smoothness} --- max jerk $\max_t \|\dddot{\mathbf{q}}(t)\|$
\end{enumerate}

%======================================================================
\newpage
\section{Experiment 1: Tabletop Pick-and-Place}
%======================================================================

\subsection{Purpose}
Validate that the joint-inertia-weighted Riemannian metric produces lower-torque, smoother paths than Euclidean planners when navigating around localized obstacles on a table surface.

\subsection{Objects Needed}

\begin{table}[h]
\centering
\begin{tabular}{@{}clllc@{}}
\toprule
\# & Object & Real Item & Size (cm) & Position wrt Base (m) \\
\midrule
1 & Table surface & Flat wooden board / MDF & $60 \times 60 \times 4$ & $[0.50,\; 0.00,\; 0.40]$ \\
2 & Tall box & Red wooden block & $16 \times 16 \times 24$ & $[0.40,\; 0.00,\; 0.55]$ \\
3 & Small box & Blue wooden block & $12 \times 12 \times 16$ & $[0.60,\; 0.15,\; 0.50]$ \\
\bottomrule
\end{tabular}
\caption{Tabletop obstacle list. Positions refer to the center of each object.}
\end{table}

\subsection{Obstacle Placement Detail}

\begin{itemize}[itemsep=2pt]
  \item \textbf{Table surface (Object 1):} A flat board elevated so the \emph{top} surface is at $z = 0.42$\,m. Use legs, risers, or stack of wood to raise it. The board spans $x \in [0.20, 0.80]$, $y \in [-0.30, 0.30]$. Clamp or weigh down so it does not shift during robot motion.
  \item \textbf{Tall red box (Object 2):} Place on the table, centered at $x=0.40$, $y=0.00$ --- directly in front of the robot. The box bottom sits on the table surface ($z=0.42$), so its center is at $z = 0.42 + 0.12 = 0.54$\,m. This is the primary obstacle.
  \item \textbf{Small blue box (Object 3):} Place on the table at $x=0.60$, $y=+0.15$ --- forward and to the left of the red box. Its center is at $z = 0.42 + 0.08 = 0.50$\,m. This creates a secondary obstacle the arm must avoid.
\end{itemize}

\subsection{Planner Obstacle Definitions}

Enter into the planner collision checker as axis-aligned boxes:
\begin{verbatim}
  Box 1: center=[0.50, 0.00, 0.40], half_extents=[0.30, 0.30, 0.02]
  Box 2: center=[0.40, 0.00, 0.55], half_extents=[0.08, 0.08, 0.12]
  Box 3: center=[0.60, 0.15, 0.50], half_extents=[0.06, 0.06, 0.08]
\end{verbatim}

\subsection{Start and Goal Configurations}

\begin{table}[h]
\centering
\begin{tabular}{@{}lllp{5cm}@{}}
\toprule
 & Joint Angles (deg) & EE Position (m) & Description \\
\midrule
Start & $[0,\; -90,\; 0,\; -90,\; 0,\; 0]$
      & $[0.00,\; 0.29,\; 1.49]$
      & Arm straight up --- standard UR10e home pose \\[4pt]
Goal  & $[68.8,\; -68.8,\; 57.3,\; -80.2,\; -90,\; 0]$
      & $[0.17,\; 0.91,\; 0.75]$
      & Arm reaching to the left-forward side of the table \\
\bottomrule
\end{tabular}
\end{table}

\subsection{Setup Procedure}

\begin{enumerate}[itemsep=2pt]
  \item Mount UR10e on a sturdy flat surface. Mark the base center with tape.
  \item Place the wooden board so its center is \SI{50}{\centi\meter} forward ($+X$) from the base, elevated on supports so the top surface is at $z = 0.42$\,m. Use a level.
  \item Place the tall red block at $x = 0.40$\,m, $y = 0.00$\,m on the board (directly in front).
  \item Place the small blue block at $x = 0.60$\,m, $y = +0.15$\,m on the board (forward-left).
  \item Measure all positions with a tape measure from the base center. Required accuracy: $\pm$\SI{2}{\centi\meter}.
  \item Teach the start and goal joint positions using the teach pendant. Verify EE positions match the table values within $\pm$\SI{3}{\centi\meter}.
  \item Dry-run at 5\% speed to confirm no collisions.
\end{enumerate}

\subsection{Experiment-Specific Metrics}
\begin{itemize}[itemsep=2pt]
  \item Total energy: $E = \sum_t |\boldsymbol{\tau}(t) \cdot \Delta\mathbf{q}(t)|$ --- expect RIT* to use 20--30\% less energy than BIT*
  \item Compare max torque on shoulder\_lift\_joint (heaviest link, 13.05\,kg)
\end{itemize}


%======================================================================
\newpage
\section{Experiment 2: Shelf Reach-In}
%======================================================================

\subsection{Purpose}
Validate that the spatially-varying metric field improves planning in constrained narrow-passage environments. The metric should inflate near shelf walls, biasing samples toward the passage center and improving clearance.

\subsection{Objects Needed}

\begin{table}[h]
\centering
\begin{tabular}{@{}clllc@{}}
\toprule
\# & Component & Real Item & Size (cm) & Position wrt Base (m) \\
\midrule
1 & Back wall & Plywood sheet & $4 \times 70 \times 60$ & $[1.22,\; 0.00,\; 0.60]$ \\
2 & Top shelf & Plywood sheet & $38 \times 70 \times 4$ & $[1.05,\; 0.00,\; 0.90]$ \\
3 & Bottom shelf & Plywood sheet & $38 \times 70 \times 4$ & $[1.05,\; 0.00,\; 0.30]$ \\
4 & Left wall & Plywood sheet & $38 \times 4 \times 60$ & $[1.05,\; +0.35,\; 0.60]$ \\
5 & Right wall & Plywood sheet & $38 \times 4 \times 60$ & $[1.05,\; -0.35,\; 0.60]$ \\
\bottomrule
\end{tabular}
\caption{Shelf components. Alternatively, use a real open-front bookshelf with matching interior dimensions.}
\end{table}

\subsection{Obstacle Placement Detail}

\textbf{Option A --- Build from plywood:}
\begin{itemize}[itemsep=2pt]
  \item Assemble a 5-sided open-front box. The \emph{front opening} faces the robot.
  \item Interior dimensions: \textbf{36\,cm deep} ($X$), \textbf{70\,cm wide} ($Y$), \textbf{60\,cm tall} ($Z$).
  \item Position the shelf so the front edge is at $x = 0.86$\,m and the back wall at $x = 1.24$\,m from the robot base.
  \item The opening center is at $y = 0.00$\,m (directly in front) and $z = 0.60$\,m (mid-height).
  \item Screw or clamp the panels together. Bolt or clamp the assembly to the floor/table so it cannot shift.
\end{itemize}

\textbf{Option B --- Use a real bookshelf/cabinet:}
\begin{itemize}[itemsep=2pt]
  \item Find any shelf with interior dims $\geq$\,35\,cm deep $\times$ 65\,cm wide $\times$ 55\,cm tall.
  \item Position with front opening facing the robot, front edge at $x \approx 0.86$\,m. Adjust planner parameters to match actual dimensions.
\end{itemize}

\subsection{Planner Obstacle Definitions}

\begin{verbatim}
  Back wall:    center=[1.22, 0.00, 0.60], half=[0.02, 0.35, 0.30]
  Top shelf:    center=[1.05, 0.00, 0.90], half=[0.19, 0.35, 0.02]
  Bottom shelf: center=[1.05, 0.00, 0.30], half=[0.19, 0.35, 0.02]
  Left wall:    center=[1.05, 0.35, 0.60], half=[0.19, 0.02, 0.30]
  Right wall:   center=[1.05,-0.35, 0.60], half=[0.19, 0.02, 0.30]
\end{verbatim}

\subsection{Start and Goal Configurations}

\begin{table}[h]
\centering
\begin{tabular}{@{}lllp{5.5cm}@{}}
\toprule
 & Joint Angles (deg) & EE Position (m) & Description \\
\midrule
Start & $[0,\; -114.6,\; 90,\; -63,\; -90,\; 0]$
      & $[0.38,\; 0.17,\; 0.86]$
      & Arm retracted, EE above and in front of the shelf opening \\[4pt]
Goal  & $[0,\; -57.3,\; 57.3,\; -90,\; -90,\; 0]$
      & $[1.02,\; 0.17,\; 0.58]$
      & Arm extended, EE inside the shelf at mid-height \\
\bottomrule
\end{tabular}
\end{table}

\subsection{Setup Procedure}

\begin{enumerate}[itemsep=2pt]
  \item Place the shelf assembly \SI{1.05}{\meter} in front of the robot base (center of shelf body). Front opening faces robot.
  \item Verify the shelf is level and stable. Bolt or clamp to floor.
  \item The shelf opening center should be at $y = 0.00$\,m and $z = 0.60$\,m.
  \item Teach start pose: arm retracted, EE at approximately $[0.38,\; 0.17,\; 0.86]$ --- above the shelf opening.
  \item Teach goal pose: arm extended into the shelf, EE at approximately $[1.02,\; 0.17,\; 0.58]$ --- inside the shelf at mid-height.
  \item Verify the goal EE position is at least \SI{5}{\centi\meter} from any wall.
  \item Dry-run at 5\% speed; observe clearance around all shelf walls.
\end{enumerate}

\subsection{Experiment-Specific Metrics}
\begin{itemize}[itemsep=2pt]
  \item \textbf{Minimum clearance:} Record full EE trajectory via FK at \SI{500}{\hertz}. Compute the min distance to any shelf wall at each timestep. RIT* should maintain higher clearance (\SI{>3}{\centi\meter}) vs.\ BIT*.
  \item \textbf{Contact detection:} If the robot touches the shelf at any point $\rightarrow$ mark trial as failure.
  \item \textbf{Number of collision checks:} Log from the planner. RIT* should need fewer due to spatially-aware sampling.
\end{itemize}


%======================================================================
\newpage
\section{Experiment 3: Overhead Reach-Across}
%======================================================================

\subsection{Purpose}
Validate that the inertia-weighted metric drives energetically efficient route selection when the direct path is globally blocked. The planner must choose to go over or under a hanging panel; the Riemannian metric should prefer the route that moves lighter joints.

\subsection{Objects Needed}

\begin{table}[h]
\centering
\begin{tabular}{@{}clllc@{}}
\toprule
\# & Object & Real Item & Size (cm) & Position wrt Base (m) \\
\midrule
1 & Hanging panel & MDF board, foam board, or cardboard & $40 \times 30 \times 5$ & $[0.25,\; 0.06,\; 0.68]$ \\
2 & Support frame & 80/20 aluminum extrusion or camera tripod crossbar & --- & Overhead, $z > 1.0$\,m \\
\bottomrule
\end{tabular}
\caption{Overhead experiment objects. Use foam board or cardboard for the panel for safety during testing.}
\end{table}

\subsection{Obstacle Placement Detail}

\begin{itemize}[itemsep=2pt]
  \item The hanging panel is suspended from an overhead frame using string or clamps. It must hang freely and \textbf{not} be supported from below.
  \item The panel center is at $[x = 0.25,\; y = 0.06,\; z = 0.68]$\,m --- in the middle of the workspace, at EE height for horizontal arm poses.
  \item The panel spans:
  \begin{itemize}
    \item $X$: from 0.05\,m to 0.45\,m
    \item $Y$: from $-0.09$\,m to 0.21\,m
    \item $Z$: from 0.655\,m to 0.705\,m
  \end{itemize}
  \item At both start and goal poses, the EE is at $z = 0.677$\,m --- \textbf{exactly at panel height}. Therefore the straight-across joint interpolation is blocked.
  \item \textbf{Safety:} Use foam board or thick cardboard for the panel rather than wood or metal. If the robot contacts it during testing, the panel deforms harmlessly.
\end{itemize}

\subsection{Planner Obstacle Definition}

\begin{verbatim}
  Box: center=[0.25, 0.06, 0.68], half_extents=[0.20, 0.15, 0.025]
\end{verbatim}

\subsection{Start and Goal Configurations}

\begin{table}[h]
\centering
\begin{tabular}{@{}lllp{5.5cm}@{}}
\toprule
 & Joint Angles (deg) & EE Position (m) & Description \\
\midrule
Start & $[-68.8,\; -90,\; 90,\; -90,\; -90,\; 0]$
      & $[0.41,\; -0.58,\; 0.68]$
      & Arm reaching to the right side of the workspace \\[4pt]
Goal  & $[68.8,\; -90,\; 90,\; -90,\; -90,\; 0]$
      & $[0.09,\; 0.71,\; 0.68]$
      & Arm reaching to the left side of the workspace \\
\bottomrule
\end{tabular}
\end{table}

The only difference between start and goal is the shoulder pan joint rotating $137.6^\circ$. The direct rotation path sends the forearm and gripper straight through the hanging panel.

\subsection{Setup Procedure}

\begin{enumerate}[itemsep=2pt]
  \item Build or position an overhead frame above the robot. An 80/20 extrusion frame, a camera rig crossbar, or even a door frame works.
  \item Hang the $40 \times 30 \times 5$\,cm panel from the crossbar using string or clamps. The panel center should be at $[0.25,\; 0.06,\; 0.68]$\,m from the robot base.
  \item Verify the panel height: with the robot in a horizontal-arm pose (shoulder lift $= -90^\circ$, elbow $= 90^\circ$), the EE should be at the same $z$ as the panel center ($z \approx 0.68$\,m).
  \item Teach start pose: EE at $[0.41,\; -0.58,\; 0.68]$ (right side). Teach goal pose: EE at $[0.09,\; 0.71,\; 0.68]$ (left side).
  \item \textbf{Verification:} Run a simple linear joint interpolation from start to goal at 5\% speed with protective stop enabled. Confirm the robot \emph{collides} with the panel --- this proves the obstacle blocks the direct path and the planner must find a detour.
  \item Ensure the support frame itself is outside the robot's reachable workspace or high enough ($z > 1.2$\,m) to not interfere.
\end{enumerate}

\subsection{Experiment-Specific Metrics}
\begin{itemize}[itemsep=2pt]
  \item \textbf{Path choice:} Record the full FK trajectory. Classify the path as ``over'' ($z_\text{max} > 0.75$\,m along the path) or ``under'' ($z_\text{min} < 0.60$\,m). RIT* should consistently prefer ``under'' because it is the lower-energy route (moves wrist joints more, shoulder less).
  \item \textbf{Shoulder torque:} Compare peak and total torque on \texttt{shoulder\_pan\_joint} and \texttt{shoulder\_lift\_joint} between planners. RIT* paths should use 20--40\% less energy since the inertia-weighted metric penalizes heavy-joint motion.
  \item \textbf{Path length in joint space:} RIT* may produce a longer Euclidean path but a shorter Riemannian path --- report both.
\end{itemize}


%======================================================================
\newpage
\section{Summary: Difficulty Progression}
%======================================================================

\begin{table}[h]
\centering
\begin{tabular}{@{}clll@{}}
\toprule
Exp & Scene & Obstacle Complexity & Why Harder Than Previous \\
\midrule
1 & Tabletop     & 2 small boxes on a surface & Baseline: obstacles are localized, most of C-space is free \\
2 & Shelf        & 5-sided enclosure       & Narrow opening constrains entry and exit \\
3 & Overhead panel & 1 large panel at EE height & Forces global path detour; tests metric-driven route selection \\
\bottomrule
\end{tabular}
\end{table}

\section{Safety Checklist}

\begin{enumerate}[itemsep=2pt]
  \item Enable UR10e protective stop mode during all tests.
  \item First run of every new path: execute at 5\% speed.
  \item Data collection runs: max \SI{0.5}{\radian\per\second} per joint.
  \item Use foam/cardboard for Experiment 3 panel.
  \item Keep emergency stop button accessible at all times.
  \item Clear the robot workspace of any objects not listed above before each experiment.
\end{enumerate}

\end{document}
